\chapter{Notas y tags}
\label{sec:man-notas}

\section{Tags}

Una etiqueta es un string corto asociado a una entidad. Tag típicos:

\begin{itemize}
    \item \code{sospechoso-c2} — entidad candidata a command-and-control.
    \item \code{cuenta-comprometida} — user con evidencia de compromiso.
    \item \code{externa-benigna} — IP/domain externo ya triado como no malicioso.
    \item \code{hub-corporativo} — entidad hub que hay que ignorar de la caza.
\end{itemize}

Crear: click derecho sobre un nodo $\to$ \emph{Tag} $\to$ escribir
el tag. Remover: click derecho $\to$ \emph{Untag}.

\section{Listar por tag}

\emph{Panel \textrightarrow{} Tags} muestra todas las etiquetas
creadas con sus conteos. Click sobre un tag lista las entidades
que lo llevan. Útil para:

\begin{itemize}
    \item Revisar todos los candidatos de una categoría.
    \item Armar el inventario de evidencia para un reporte.
\end{itemize}

\section{Notas}

Una nota es texto libre asociado a una entidad, un path
específico, o un hunt. Se crean desde:

\begin{itemize}
    \item Click derecho sobre un nodo $\to$ \emph{Note}.
    \item Sobre un path en los resultados de hunt $\to$
    \emph{Annotate}.
    \item Sobre un hunt en el historial $\to$ \emph{Note}.
\end{itemize}

Soporta Markdown básico. Quedan persistidas en la sesión.

\section{Workflow de reporte}

Flujo típico de construir un reporte desde dentro de la app:

\begin{enumerate}
    \item Ingerir datos.
    \item Correr hunts relevantes; tagear paths sospechosos.
    \item Explorar los paths top y tagear entidades clave con
    \code{sospechoso-*} o \code{ioc}.
    \item Agregar notas con la narrativa en cada entidad clave
    (``IP externa con beacon a intervalos regulares, ver captura
    \#12'').
    \item Al final: \emph{Listar por tag \textrightarrow{} ioc}
    da el set de entidades para exportar como IOCs.
    \item \emph{Historial de hunts} listado con sus notas da la
    narrativa completa.
\end{enumerate}

\section{Interacción con DSL}

Una hipótesis puede filtrar por tag:

\begin{lstlisting}[basicstyle=\ttfamily\small]
IP {tag("externa-benigna")} -[Connect]-> Host
\end{lstlisting}

O lo más común: \textbf{excluir} entidades ya triadas como no
maliciosas para que no contaminen resultados de nuevos hunts.

Los tags viven en la sesión. No se propagan entre sesiones a menos
que las exportes explícitamente.
